\section{NAME}

\textbf{codata }- Command line for codata

\section{SYNOPSIS}

\begin{verbatim}
codata [OPTIONS] [REGEX_PATTERN ... ]
\end{verbatim}

\section{DESCRIPTION}

\textbf{codata is a command line interface which prints all the codata}
constants.

The current values are from 2022. Older values can be retrieved if
needed and the output can be filtered with REGEX PATTERNS.

\section{OPTIONS}

\begin{description}
\item[\textbf{\textbf{-\/-year, -y YEAR}}]
Year of the \textbf{codata constants:} 2022, 2018, 2014, 2010
\item[\textbf{\textbf{-\/-pattern, -p PATTERN}}]
Regex pattern for filtering the constants.
\item[\textbf{\textbf{-\/-value, -a}}]
Show only the value.
\item[\textbf{\textbf{-\/-error, -e}}]
Show only the uncertainty.
\item[\textbf{\textbf{-\/-usage}}]
Show usage text and exit.
\item[\textbf{\textbf{-\/-help}}]
Show help text and exit.
\item[\textbf{\textbf{-\/-verbose}}]
Display additional information.
\item[\textbf{\textbf{-\/-version}}]
Show version information and exit.
\end{description}

\section{NOTES}

You may replace the default options from a file if your first options
begin with @file. Initial options will then be read from the "response
file" "file.rsp" in the current directory.

If "file" does not exist or cannot be read, then an error occurs and the
program stops. Each line of the file is prefixed with "options" and
interpreted as a separate argument. The file itself may not contain
@file arguments. That is, it is not processed recursively.

For more information on response files see
https://urbanjost.github.io/M\_CLI2/set\_args.3m\_cli2.html

\section{EXAMPLE}

Minimal example

\begin{verbatim}
     codata
     codata -y 2018 molar electron
     codata -y 2014 -p 'molar.*gas','electron.*eV'
     codata '[B,b]oltzmann.*eV'
\end{verbatim}

\section{SEE ALSO}

\textbf{codata(3)}
